\documentclass[12pt,a4paper]{article}
\usepackage[margin=2cm]{geometry}
\usepackage[spanish]{babel}
\usepackage{graphicx, multicol, latexsym, amsmath, amssymb}
\usepackage{array}
\usepackage{tabularx}
\usepackage{setspace}
\usepackage{titlesec}
\usepackage{subcaption}
\usepackage{float}
\usepackage{url}
\usepackage{hyperref}

% Ajuste de espacio entre secciones
\titlespacing*{\section}{0pt}{0.5em}{0.5em}

\renewcommand{\arraystretch}{1.3}

\begin{document}

% ---------------------------
% Encabezado con datos
% ---------------------------
\noindent
\begin{tabularx}{\textwidth}{|X|X|}
\hline
\textbf{Carné:} 2020045294 & \textbf{Nombre:} Justin Jaffeth Corea Masís \\
\hline
\textbf{Correo:} coreajustin288@estudiantec.cr & \textbf{Semana:} 6 del 10/09/2025 al 17/09/2025 \\
\hline
\end{tabularx}

\vspace{0.5cm}

% ---------------------------
% Tabla de actividades
% ---------------------------
\section*{Actividades realizadas}

% \begin{tabularx}{\textwidth}{|>{\raggedright\arraybackslash}p{8cm}|c|c|}
\begin{tabularx}{\textwidth}{|>{\raggedright\arraybackslash}p{12cm}|c|c|}
\hline
\textbf{Actividad} & \textbf{Fecha} & \textbf{Horas} \\
\hline
	Cambié redacción de introducción de tesis & 11/09/2025 & 2 \\
\hline
	Reunión de aspectos operativos de TFG con profesor Miguel & 12/09/2025 & 2 \\
\hline
	Inicié marco teórico y modifiqué capítulo de solución en tesis & 14/09/2025 & 2 \\
\hline
	Hice más experimentos con el OCR & 15/09/2025 & 5 \\
\hline
	Continué experimentando con el OCR & 16/09/2025 & 2 \\
\hline
	Continué redacción de marco teórico & 16/09/2025 & 2 \\
\hline
\end{tabularx}

\vspace{1cm}

% ---------------------------
% Firma
% ---------------------------
\noindent
\textbf{Firma del Profesor Asesor:}

\vspace{2cm}

\noindent\rule{8cm}{0.4pt}  
Nombre y firma

\newpage

\section{Reconocimiento de placas OCR}
Generé más imágenes sintéticas de los tipos de placa subrepresentadas.
Busque las placas que estaban dando mal reconocimiento en los mejores modelos.
Eliminé las que me parecía muy complejas de reconocer incluso a ojo humano y 
reentrené. Además, corregí unas imágenes mal etiquetadas
en el conjunto de imágenes reales.
El modelo mejoró solamente un 1\% respecto al baseline,
pero algunas siguen dando resultados incorrectos a pesar de
no ser realmente difíciles de reconocer (en mi opinión).

\begin{figure}[H]
	\centering
	\begin{subfigure}[b]{0.3\textwidth}
		\includegraphics[width=\textwidth]{/home/akai/Pictures/fallo-4.png}
	\end{subfigure}
	\hfill
	\begin{subfigure}[b]{0.3\textwidth}
		\includegraphics[width=\textwidth]{/home/akai/Pictures/fallo-5.png}
	\end{subfigure}
	\hfill
	\begin{subfigure}[b]{0.3\textwidth}
		\includegraphics[width=\textwidth]{/home/akai/Pictures/fallo-6.png}
	\end{subfigure}

	\begin{subfigure}[b]{0.3\textwidth}
		\includegraphics[width=\textwidth]{/home/akai/Pictures/fallo-7.png}
	\end{subfigure}
	\hfill
	\begin{subfigure}[b]{0.3\textwidth}
		\includegraphics[width=\textwidth]{/home/akai/Pictures/fallo-2.png}
	\end{subfigure}
	\hfill
	\begin{subfigure}[b]{0.3\textwidth}
		\includegraphics[width=\textwidth]{/home/akai/Pictures/fallo-3.png}
	\end{subfigure}
	\caption{Ejemplos de imágenes de validación en la que todos los modelos fallan}
\end{figure}

\end{document}
